% !TEX root = ../algebra_02.tex

\section{Расширение, порождённое данным конечным множеством}

Пусть $L/K$ --- расширение, и $P \subset L$ --- некоторое подмножество.
Определим $K(P)$ как пересечение всех промежуточных полей $F$, содержащих $K$ и $P$:
\[ K(P) := \bigcap_{\substack{K \subset F \subset L \\ F \text{ --- поле} \\ P \subset F}} F \]
Очевидно, что $K(P)$ является полем. Оно называется полем, полученным из $K$ \textbf{присоединением элементов} из множества $P$.

Если множество элементов конечно ($P = \{x_1, \dots, x_n\}$), используют обозначение $K(x_1, \dots, x_n)$.

\begin{Definition}{Конечно порождённое расширение}{}
	Расширение $L/K$ называется \textbf{конечно порождённым}, если существуют элементы $x_1, \dots, x_n \in L$ такие, что $L = K(x_1, \dots, x_n)$.
\end{Definition}

\begin{Definition}{Простое расширение}{}
	Расширение $L/K$ называется \textbf{простым}, если оно порождается одним элементом: $L = K(x)$ для некоторого $x \in L$.
\end{Definition}

\begin{Remark}{}{}
	Если элемент $a$ трансцендентен над $K$, то расширение $L = K(a)$ называется простым трансцендентным.
	Если $a$ алгебраичен, то $L = K(a)$ называется простым алгебраическим расширением.
\end{Remark}

\begin{Definition}{Алгебраическое расширение}{}
	Расширение $L/K$ называется \textbf{алгебраическим}, если все элементы поля $L$ алгебраичны над $K$. В противном случае расширение $L/K$ называется \textbf{трансцендентным}.
\end{Definition}
